\section*{Chapter \ref{ch:rc:margin-models} --- Margin Models}

\begin{solution}{exo:rc:margin-models:1}
$4\%\times100 = 4$ million for the ten-year and $1\%\times100 = 1$ million for the one-year: USD~5 million gross.
\end{solution}

\begin{solution}{exo:rc:margin-models:2}
Its 99\% quantile is the third-worst of about 250 five-day moves: a few new volatile days add a few new
tail observations but do not move the quantile much until they outnumber the old tail, and the calm days
of the past year still dominate the window.
\end{solution}

\begin{solution}{exo:rc:margin-models:3}
$\mathrm{NGR} = 30/120 = 0.25$; net margin $= (0.4+0.6\times0.25)\times\text{gross} = 0.55\times\text{gross}$: a 45\% reduction.
\end{solution}

\begin{solution}{exo:rc:margin-models:4}
SIMM's risk weights are calibrated on long histories including stress periods and change only at
recalibration; they do not react to current volatility. Clearing-house models track volatility (filtered
or short look-back), so their margin rose with it.
\end{solution}

\begin{solution}{exo:rc:margin-models:5}
The peak is set by the model margin at the height of the stress, which the tools do not change (the
buffer is released, the stressed weight is below the March level); they only raise the margin held
before, so the increase is smaller.
\end{solution}

\begin{solution}{exo:rc:margin-models:6}
The time and price impact of liquidating a position that is large relative to daily volume: the close-out
takes longer than the margin period of risk and moves prices, which a model of normal-size moves over a
fixed horizon does not see.
\end{solution}

\begin{solution}{exo:rc:margin-models:7}
Expected $2\,177\times1\% = 21.8$; Kupiec likelihood ratio 6.86, p-value 0.009. Overlapping five-day windows make
consecutive days share four of their five days of moves, so exceptions come in runs; the test assumes
independence and overstates the evidence.
\end{solution}

\begin{solution}{exo:rc:margin-models:8}
Coverage must be tested out of sample and over stress: a model calibrated to calm months covers calm
months and fails when volatility jumps, as March 2020 showed. The test is the next month, not the last.
\end{solution}

\begin{solution}{pb:rc:margin-models:1}
\textbf{1.} USD~4.21 million in February; 10.30 million on 23 March.
\textbf{2.} USD~6.09 million.
\textbf{3.} Volatility: the positions are fixed in the replay.
\textbf{4.} USD~1.06 million (from 4.19 to 5.25).
\textbf{5.} It covers the current risk: the historical model under-margined the book during the stress,
which is when a default is most likely.
\textbf{6.} USD~5.26 million held; an increase of 5.04 million.
\textbf{7.} USD~4.79 million held; an increase of 5.51 million.
\textbf{8.} The funding of the extra margin held in calm years (about 1 million under the buffer, 0.6 under
the stressed weight in February), at the member's funding spread over what margin earns.
\textbf{9.} After volatility has normalised, and gradually: rebuilding it at once would be another margin
call just as markets recover.
\textbf{10.} From a 25\% increase over five days to up to 80\%.
\textbf{11.} Ten-day historical USD~4.15 million; sensitivity-based 8.98 million; schedule 13.19 million
gross, 9.23 million net.
\textbf{12.} Netting across all its cleared trades, lower margin than the schedule, and no bilateral
counterparty risk; in exchange for exposure to the clearing house's procyclical model.
\textbf{13.} Every figure is below EUR~50 million: under the threshold, no initial margin would be
exchanged bilaterally for this book alone.
\textbf{14.} In a crisis both rise: variation margin moves daily with prices and must be paid in cash;
initial margin rises with volatility. Together they drain liquidity at the worst time.
\textbf{15.} The clearing house's own model (or a replica) run on stress scenarios, not the calm-period margin.
\textbf{16.} A feature for the clearing house's safety, a flaw for the system's liquidity: it transfers
stress to members when they can least bear it.
\textbf{17.} Members, through higher margin in calm times, which they must fund.
\textbf{18.} By the margin increase in severe historical and hypothetical stresses (here more than 6 million
in three weeks), plus variation margin, with collateral that can be delivered on the day.
\textbf{19.} Named result: \emph{the margin calls of March 2020}: the filtered model called USD~6.09
million of new margin; with a 25\% buffer 5.04 million and with a 25\% stressed weight 5.51 million.
\textbf{20.} Protection of the counterparty against predictability and liquidity for the payer.
\end{solution}

\begin{solution}{iq:rc:margin-models:1}
It covers the loss of closing out a defaulter's portfolio over the margin period of risk at a high
confidence. Clearing houses use historical, filtered or scenario models on the member's whole portfolio,
at least 99\% over a few days.
\iqlookfor{purpose, horizon, confidence and method.}
\end{solution}

\begin{solution}{iq:rc:margin-models:2}
Margin rises with volatility, so members must post more exactly in stress, which forces sales and adds
to stress. Tools: buffers released in stress, stressed-period weights, long look-backs and volatility
floors.
\iqlookfor{the feedback and three tools.}
\end{solution}

\begin{solution}{iq:rc:margin-models:3}
Compute sensitivities per \omterm{def:rc:market-risk-measures:factor}{risk factor} (delta by tenor, vega, curvature), multiply by calibrated risk
weights, aggregate within buckets with correlations and across buckets and classes; parameters are
recalibrated periodically on stress-inclusive histories.
\iqlookfor{sensitivities, weights, correlations and calibration.}
\end{solution}

\begin{solution}{iq:rc:margin-models:4}
Different models (historical, filtered, scenario), look-backs, margin periods of risk, confidence levels,
APC tools, add-ons for concentration and liquidity, and \omterm{def:rc:counterparty-exposure:netting}{netting sets}.
\iqlookfor{model and parameter differences.}
\end{solution}

\begin{solution}{iq:rc:margin-models:5}
Compare realised losses over the margin period of risk with the margin held, per portfolio, and test
frequency and clustering; hard because portfolios change, windows overlap, stress is rare, and the
right loss includes liquidation costs.
\iqlookfor{the procedure and its difficulties.}
\end{solution}

\begin{solution}{iq:rc:margin-models:6}
Margin calls, initial and variation, were large, fast and synchronised; firms that had not forecast them
sold assets into falling markets. Margin models need to be predictable, and firms need liquidity plans
built on stressed margin.
\iqlookfor{liquidity effects and predictability.}
\end{solution}
